% If you intend to adapt the material in this textbook
% in any way, then you must set the \adaptername and
% \adapteremail commands on lines 13 and 14 of this file
% to your full name and email address, respectively.
%
% Example usage:
% \newcommand{\adaptername}{Exampleface McLastname}
% \newcommand{\adapteremail}{exampleface-mclastname@example.web}
%
% Any duplications, derivations or adaptations of this work
% must be released under the Creative Commons BY-SA 4.0 licence.

\newcommand{\adaptername}{} % Your full name
\newcommand{\adapteremail}{} % Your email address

% Language
% 0:en_GB 1:en_US
\newcounter{englishvariant}\setcounter{englishvariant}{0}

% Page layout
% 0:7"x10", 1:a4, 2:usletter, 3:tablet, 4:phone
\newcounter{bookformat}\setcounter{bookformat}{0}

% Book version - if you are adapting the book, please leave
% these unchanged so that the original version can be traced.
\newcounter{versionmajor}\setcounter{versionmajor}{1}
\newcounter{versionminor}\setcounter{versionminor}{0}
\newcounter{versionpatch}\setcounter{versionpatch}{-1}
\newcounter{versionstable}\setcounter{versionstable}{0}

\documentclass[10pt]{book}

% Includes
\input{book/includes/_includes.tex}

% Uncomment to remove hyperlinks (useful for print versions)
% \hypersetup{draft}

\begin{document}

% Uncomment to display cover spread
% \includepdf[pages=-,fitpaper=true,noautoscale=false]{book/media/spread.pdf}

\frontmatter

% Title page
\newgeometry{centering,bindingoffset={0in}}
% !TeX root = ../../infdesc.tex

\begin{center}
\thispagestyle{empty}

{\Huge \mdseries \fontfamily{bch}\selectfont
An Infinite Descent into\\
Pure Mathematics}

{\Large \mdseries
\textit{Solution Book}}

\vfill
\includegraphics{book/media/logo.pdf}

\vfill

\textsc{by Fisher E. Yu}\\
\ifadapted{\textit{with adaptations by \adaptername}}

\vfill

\iforiginal{%
\textit{Version \arabic{versionmajor}.\arabic{versionminor}%
\ifnumgreater{\value{versionpatch}}{0}{, revision \arabic{versionpatch}}{}%
\ifnumless{\value{versionpatch}}{0}{ preview}{}}%
\ifnumequal{\value{versionstable}}{0}{ {\color{red}(unstable build)}}{}\\
}
\textit{Last updated on \customdate\today}%
\iforiginal{\\
\url{\bookurl}}
\end{center}
\restoregeometry

% Table of contents
\bookmark[page=1,level=0]{Contents}
\tableofcontents

\mainmatter

\setcounter{chapter}{-1}

\input{book/getting-started/_getting-started.tex}

\part{Core concepts}
\label{ptCoreConcepts}

\input{book/logical-structure/_logical-structure.tex}
\input{book/sets/_sets.tex}
\input{book/functions/_functions.tex}
\input{book/mathematical-induction/_mathematical-induction.tex}
\input{book/relations/_relations.tex}
\input{book/finite-infinite-sets/_finite-infinite-sets.tex}

\part{Topics in pure mathematics}
\label{ptTopics}

\input{book/number-theory/_number-theory.tex}
\input{book/enumerative-combinatorics/_enumerative-combinatorics.tex}
\input{book/real-numbers/_real-numbers.tex}
\input{book/infinity/_infinity.tex}
\input{book/probability-theory/_probability-theory.tex}
\input{book/additional-topics/_additional-topics.tex}

\appendix

\begin{appendices}

\renewcommand{\sectionmark}[1]{\markboth{\leftmark}{Section \thesection.\ #1}}
\renewcommand{\chaptermark}[1]{\markboth{Appendix \thechapter.\ #1}{\rightmark}}

\part*{Appendices}
\addcontentsline{toc}{part}{Appendices}

\input{book/proof-writing/_proof-writing.tex}
\input{book/miscellany/_miscellany.tex}

\end{appendices}

\input{book/licence/_licence.tex}

\cleardoublepage

\end{document}